\section{Finite counterfactuals and potential outcomes: proofs}
\label{sec:annex-p16-finite-counterfactuals}

\begin{annexitem}{annex:p16.counterfactual-laws}{thm:p16-counterfactual-laws}{The identities follow from the actual replaced equations}
\begin{proof}
All modified systems are finite and acyclic, with total functions, so
\cref{thm:p16-causal-solution} supplies unique solutions at every
outcome. At an intervened coordinate its equation is the prescribed
constant, proving effectiveness. If the factual solution already has
\(V_I(\omega)=x\), that tuple satisfies every retained equation
and each replacement equation. Uniqueness proves factual consistency.

Now let \(v=V[I=x](\omega)\) and suppose \(v_J=y\).
The tuple \(v\) satisfies every equation of the already modified
model. Replacing its equations at \(J\) by the already attained
values \(y\) leaves \(v\) a solution: equations outside \(J\)
are unchanged, and those inside \(J\) are satisfied by the supposition.
Uniqueness gives the composition conclusion, including when \(I,J\)
overlap. On that overlap the attained values necessarily agree with
the constants previously imposed at \(I\).

For arbitrary successive assignments, examine the final equation at
each vertex. On \(J\) it is the constant \(y_j\), on
\(I\setminus J\) it is the constant \(x_i\), and elsewhere
it is the original equation. This is exactly the stated single
replacement on the union. Reversing the order gives the same equations
if the two assignments agree on the intersection. If they disagree,
the later assignment supplies the constants there; no commutation
claim is made in that case.

For an explicit failure of substituting another world's value, take
the singleton-outcome model \(A=0,B=A,Y=B\) with graph
\(A\to B\to Y\). In the world \(A=0\), both \(B\) and
\(Y\) are zero. The world \(A=1\) instead gives \(B=1\).
Replacing \(B\) by that latter value while keeping \(A=0\)
gives \(Y=1\), not the response in the \(A=0\) world. All
equations are total and acyclic, so the failed equality cannot be
repaired by appealing to solution ambiguity.

Finally put \(H=N\setminus(I\cup\operatorname{de}_G(I))\).
Every parent of a vertex in \(H\) is in \(H\): a parent in
\(I\) or reachable from \(I\) would give a directed path to
the child after appending its edge. All equations in \(H\) are
retained. Induction in an original topological order restricted to
\(H\) shows that their inputs and outputs agree before and after
replacement. This proves the last assertion. None of these arguments
uses probabilities or disturbance independence.
\end{proof}
\end{annexitem}

\begin{annexitem}{annex:p16.counterfactual-calculation}{thm:p16-counterfactual-calculation}{The posterior pushforward and its finite calculation}
\begin{proof}
For every possible tuple value \(v\) of \(Q\), form its fiber
\(E_v=\{\omega:Q(\omega)=v\}\). The fibers partition
\(\Omega\), so their masses are nonnegative and sum to one.
Summing over all tuple values with the specified coordinate in one
world gives exactly that world's solution fiber, proving the marginal
claim. For positive evidence \(E\), finite conditioning gives
\[
 \mathbb P(Q\in D\mid E)
  =\frac{\sum_{\omega\in E,\ Q(\omega)\in D}q(\omega)}{\mathbb P(E)}
  =\sum_{\omega:Q(\omega)\in D}q_E(\omega).
\]
The weights \(q_E\) sum to one by the same conditioning formula.
Computing them, evaluating each of the finitely many world systems
in topological order at each of the finitely many outcomes, and
summing the indicated weights therefore terminates and gives exactly
this probability. There are at most
\(|\Omega|\,|\mathcal W|\,|N|\) structural-coordinate evaluations
once the orders and posterior weights are available.

For any fixed disturbance tuple \(u\in\mathcal U\), evaluate
the replaced equations using its component \(u_i\) at vertex \(i\).
The total functions and acyclic order define a unique solution even
when that particular tuple has probability zero. Induction in the
order shows that two original outcomes with the same tuple \(U=u\)
give that same solution in every world. Denote the resulting queried
tuple by \(\widehat Q(u)\). Grouping the preceding sum by \(u\)
gives
\[
 \mathbb P(Q\in D\mid E)
   =\sum_{u:\widehat Q(u)\in D}r_E(u),
 \qquad
 r_E(u)=\sum_{\omega:U(\omega)=u}q_E(\omega).
\]
The event \(E\) may contain information beyond the coordinates of
\(U\); the grouped posterior still gives the exact pushforward
because all queried worlds depend only on \(U\).

To see why the joint posterior matters, take two independent fair
disturbance bits and evidence that they agree. This event has
probability \(1/2\). Its posterior has masses \(1/2,1/2\) on
\((0,0),(1,1)\) and zero on the other pairs. Both posterior marginals
are fair, but their product would assign \(1/4\) to each pair.
Thus conditioning has destroyed their independence. If factual
observability of this evidence is desired, add endogenous roots equal
to the two bits and a child that indicates their equality, with a
constant disturbance; observing that child as one is precisely this
event. Prediction must use the computed joint posterior.
\end{proof}
\end{annexitem}

\begin{annexitem}{annex:p16.finite-twin-construction}{thm:p16-finite-twin-construction}{Duplicated equations with one shared background root}
\begin{proof}
The new graph has one shared root and \(|N|\,|\mathcal W|\)
copied vertices. Put the root first, and then for each world list its
vertices in an original topological order. Each within-world parent
precedes its child, and the root precedes every vertex that uses it.
There is no edge from a copied vertex back to the root or to another
world. This is a topological order, proving acyclicity.

On the original probability space set the root's disturbance equal
to the full tuple \(U\); all other disturbances are constant.
This new disturbance family is mutually independent: there is just
one possibly nonconstant member, and any tuple probability is its
probability multiplied by the indicators that all constant coordinates
have their required values. This does not require the components of
\(U\) to be independent. The root's value is \(U(\omega)\).
Induct within world \(k\). An intervened copied vertex gives
\(x_i^k\), exactly as the original intervention does. At any other
vertex, the preceding copied parent values agree with their original
world values by induction, and the root supplies the same disturbance
component \(U_i(\omega)\). Applying the same total function
gives the same output. This proves every stated pointwise identity.

For fixed \(u\), the product of indicators in the theorem is one
exactly when every copied equation is satisfied by the proposed full
assignment. Uniqueness makes that equivalent to its being the evaluated
tuple at background \(u\). Multiplication by \(r(u)\) and summation
therefore gives the full joint mass. Marginalizing the unqueried
coordinates gives the law of \(Q\). Conditioning on any retained
positive event \(E\subseteq\Omega\) preserves the pointwise
identities; equivalently, replace \(r\) by \(r_E\) in the formula.
For factual coordinate evidence, include the factual world with empty
intervention; its copied evidence event is exactly the original one.

For contrast, let one original endogenous variable equal a fair
disturbance bit, and take two worlds with empty interventions. Their
two copied values are equal pointwise, so the joint zero-zero
probability is \(1/2\). Redrawing an independent bit for each
world would give \(1/4\) there. Thus the single-world marginals
do not justify replacing the shared background by independent copies.
\end{proof}
\end{annexitem}

\begin{annexitem}{annex:p16.binary-po-bridge}{thm:p16-binary-po-bridge}{Consistency and an explicit converse representation}
\begin{proof}
At every outcome the binary variable \(A\) has one of the values
zero and one. Factual consistency in
\cref{thm:p16-counterfactual-laws} gives \(Y=Y(0)\) in the
first case and \(Y=Y(1)\) in the second. This is exactly
\(Y=(1-A)Y(0)+AY(1)\). The marginal identity is immediate from
the definition of \(Y(a)\) as the response in the model with the
treatment equation replaced. The covariate tuple is unchanged because
its vertices are nonintervened nondescendants.

For the converse, take an abstract finite probability space
\((\Omega_0,q_0)\) carrying functions
\(a_0(\omega),\ell_0(\omega),y_0(\omega),y_1(\omega)\)
for its treatment, covariates and potential responses. Introduce a
latent root \(R\) with value set \(\Omega_0\) and distribution
\(q_0\), and endogenous variables defined by
\[
 L=\ell_0(R),\qquad A=a_0(R),\qquad
 Y=(1-A)y_0(R)+Ay_1(R).
\]
Choose for \(Y\) the finite union of the ranges of \(y_0,y_1\),
and for \(L\) its finite range. The allowed edges are
\(R\to L\), \(R\to A\), \(R\to Y\), and \(A\to Y\).
Use one disturbance of law \(q_0\) at the root and constant
disturbances elsewhere. The graph is acyclic in order \(R,L,A,Y\)
and all its functions are total on their indicated finite input sets.

At root outcome \(\omega\), factual evaluation gives the original
\(A,L\), and gives the original \(Y\) by the assumed consistency.
Replacing \(A\)'s equation by \(a\) gives exactly \(y_a(\omega)\),
while \(L\) remains fixed. Hence the entire joint law, including
both potential outcomes, is represented. This construction represents
an already specified finite law; it does not discover the latent root
or its distribution from observed data. The shared root can carry
dependence between treatment and potential responses, so the construction
does not assert exchangeability.
\end{proof}
\end{annexitem}

\begin{annexitem}{annex:p16.back-door-exchangeability}{thm:p16-back-door-exchangeability}{The structural justification concerns potential outcomes}
\begin{proof}
For a fixed \(a\), use the auxiliary model constructed in the
proof of \cref{thm:p16-finite-back-door-adjustment}. It retains
the original treatment equation, replaces every treatment input at
its children by the fixed value \(a\), and deletes outgoing
treatment edges. The complete pointwise comparison in that proof
identifies its treatment coordinate as the original \(A\), its
response as \(Y(a)\), and its covariate tuple as the original
\(L\), since the covariates are nondescendants of treatment.

Every path from its treatment vertex starts with an incoming arrow,
since outgoing treatment edges were deleted. If it were active given
the covariates, its collider-activation paths would also exist in
the original graph, giving an active original back-door path. The
criterion excludes this. Thus treatment and response are d-separated
in the auxiliary graph given the covariates. Its disturbances retain
mutual independence, so \cref{thm:p16-finite-global-markov} gives
\(Y(a)\mathrel{\perp\!\!\!\perp}A\mid L\).
This is the conditional-independence conclusion before any division
by a treatment-stratum probability in that proof. It therefore holds
without an overlap assumption and for each fixed \(a\), including
treatment values absent observationally. No assertion that such a
graph or independence hypothesis is determined by the observed law
enters the argument.
\end{proof}
\end{annexitem}

\begin{annexitem}{annex:p16.po-ate-identification}{thm:p16-po-ate-identification}{Marginal exchangeability and the population weighting}
\begin{proof}
Fix \(a\in\{0,1\}\) and a stratum \(\ell\) with positive
probability. Exchangeability gives
\[
 \mathbb P(Y(a)=y,A=a\mid L=\ell)
 =\mathbb P(Y(a)=y\mid L=\ell)\mathbb P(A=a\mid L=\ell).
\]
The last factor is positive by the overlap assumption. Dividing
therefore gives
\[
 \mathbb P(Y(a)=y\mid L=\ell)
 =\mathbb P(Y(a)=y\mid A=a,L=\ell)
 =\mathbb P(Y=y\mid A=a,L=\ell),
\]
where consistency proves the last equality on the conditioning event.
Finite total probability over the positive strata proves the stated
marginal formula. Multiply that equality by each finite response
value and sum to obtain
\(\mathbb E[Y(a)]=\sum_\ell\mu_a(\ell)\mathbb P(L=\ell)\).
Subtract the formula for zero from the formula for one. Finite
linearity gives \(\mathbb E[Y(1)-Y(0)]\), proving the ATE expression.

Every term uses only the observed joint law of \((A,L,Y)\), with
the specified positive denominators. Any two models satisfying these
assumptions and sharing that observed law have identical terms and
hence identical displayed targets. Null strata contribute no mass
and require no conditional values.

The proof used exchangeability separately for each potential outcome.
To see that this does not imply joint exchangeability, let
\(Y(0),Y(1)\) be independent fair bits, put
\(A=Y(0)\oplus Y(1)\), take \(L\) constant, and define
\(Y=Y(A)\). The four equally likely potential pairs give fair
\(A\), and either individual potential outcome has one zero and
one one in each treatment stratum, so each is independent of \(A\).
The pair, however, determines \(A\). For example,
\[
 \mathbb P(Y(0)=0,Y(1)=0,A=0)=1/4
 \ne \mathbb P(Y(0)=0,Y(1)=0)\mathbb P(A=0)=1/8.
\]
Consistency and overlap both hold in this example. Thus joint
potential-outcome exchangeability was not a hidden premise. The
separate coupling theorem below shows explicitly why the identified
marginals need not identify the joint potential-outcome law.
\end{proof}
\end{annexitem}

\begin{annexitem}{annex:p16.po-att-identification}{thm:p16-po-att-identification}{Only the missing control response needs exchangeability}
\begin{proof}
Take a stratum with \(\mathbb P(A=1,L=\ell)>0\). Its treated
conditional response is defined, and the positivity hypothesis also
defines its control conditional response. Exchangeability of \(Y(0)\)
and \(A\) given \(L\), divided by each of the positive treatment
probabilities within this stratum, gives
\[
 \mathbb E[Y(0)\mid A=1,L=\ell]
 =\mathbb E[Y(0)\mid L=\ell]
 =\mathbb E[Y(0)\mid A=0,L=\ell]
 =\mu_0(\ell).
\]
The final equality uses consistency on \(A=0\). On the treated
event, consistency directly gives
\(\mathbb E[Y(1)\mid A=1,L=\ell]=\mu_1(\ell)\), without
an exchangeability assumption for \(Y(1)\). Subtract these
equalities and use conditional finite linearity. Then partition the
positive event \(A=1\) by its covariate strata to obtain
\[
 \mathbb E[\Delta\mid A=1]
 =\sum_{\ell:\mathbb P(A=1,L=\ell)>0}
    [\mu_1(\ell)-\mu_0(\ell)]
       \frac{\mathbb P(A=1,L=\ell)}{\mathbb P(A=1)}.
\]
This is the claimed formula. Its ratios and means depend only on the
observed law. Equality of that law across two models in the class
therefore gives equality of ATT. Strata with no treated mass have
weight zero in this target and were never used for division.
\end{proof}
\end{annexitem}

\begin{annexitem}{annex:p16.cf-worked-calculation}{thm:p16-cf-worked-calculation}{The full table, posterior, and two effect weightings}
\begin{proof}
The eight outcome weights \(1/8\) form the product of a fair
\(U_L\) law and a uniform four-valued \(U_A\) law. Adding the
constant response disturbance preserves mutual independence. The
topological order \(L,A,Y\) evaluates the stated total equations.
At \(L=0\), only \(U_A=1\) is treated, giving \((0,1,1)\)
mass \(1/8\); the three other outcomes give \((0,0,0)\) mass
\(3/8\). At \(L=1\), the values \(U_A=1,2,3\) are treated,
giving \((1,1,3)\) mass \(3/8\); the last value gives
\((1,0,1)\) mass \(1/8\). This proves the complete body table.

The two potential responses at each outcome are
\[
 Y(0)=L,\qquad Y(1)=1+2L,
 \qquad\Delta=1+L.
\]
Within a fixed \(L\) stratum each potential outcome is constant,
so its joint conditional probability with any treatment event is
its zero-or-one conditional probability times that treatment
probability. This proves both exchangeability assumptions. All four
treatment-covariate combinations in the table have positive mass.
Their observed means are
\[
\begin{array}{c|cc|cc}
 \ell&\mu_0(\ell)&\mu_1(\ell)&\mathbb P(L=\ell)&
                         \mathbb P(L=\ell\mid A=1)\\\hline
 0&0&1&1/2&1/4\\
 1&1&3&1/2&3/4
\end{array}
\]
The treated weights follow by dividing the masses \(1/8,3/8\)
by \(\mathbb P(A=1)=1/2\). Consequently
\[
 \operatorname{ATE}=\tfrac12(1)+\tfrac12(2)=\tfrac32,
 \qquad
 \operatorname{ATT}=\tfrac14(1)+\tfrac34(2)=\tfrac74.
\]
The observed treated response mean is
\(\frac14(1)+\frac34(3)=5/2\), while the control mean is
\(\frac34(0)+\frac14(1)=1/4\). Their difference is \(9/4\).

For abduction, the factual event \(A=1,Y=3\) consists exactly
of the three original outcomes \((U_L,U_A)=(1,1),(1,2),(1,3)\).
It has mass \(3/8\), and each compatible posterior weight is
\((1/8)/(3/8)=1/3\). For action, replace the equation for
\(A\) by zero while retaining these outcomes and weights. For
prediction, the unchanged response equation gives \(Y=L=1\)
at all three. The posterior probability of \(Y(0)=1\) is thus
\(1/3+1/3+1/3=1\).

Without that evidence, the same intervention gives \(Y(0)=L\)
with success probability \(1/2\). Among observed controls, only
the outcome \((1,4)\) has response one, and its mass \(1/8\)
divided by the control mass \(1/2\) gives \(1/4\). Every
conditional probability in these calculations has an explicit
positive denominator. The shared-background twin representation
gives the same results by its proved pointwise identities.
\end{proof}
\end{annexitem}

\begin{annexitem}{annex:p16.po-coupling-nonidentification}{thm:p16-po-coupling-nonidentification}{All single-world experiments leave the coupling undetermined}
\begin{proof}
Let \(\kappa\) be any probability mass on pairs
\((i,j)\in\{0,1\}^2\). On the eight-point space of triples
\((a,i,j)\), put mass \(\frac12\kappa(i,j)\). Define
\(U_A=a\), \(U_Y=(i,j)\), and the endogenous equations
\[
 A=U_A,\qquad
 Y=f_Y(A,(i,j)),\qquad
 f_Y(0,(i,j))=i,\quad f_Y(1,(i,j))=j.
\]
The disturbances are independent because their joint mass is the
product of a fair treatment mass and \(\kappa\). Both functions
are total on their full finite domains, and the allowed DAG is
\(A\to Y\). The two potential outcomes are exactly
\(Y(0)=i,Y(1)=j\). Thus treatment is independent of their
entire pair, with either treatment value of positive probability.

Choose \(\kappa\) first to put mass \(1/2\) at each of
\((0,0),(1,1)\), and then to put mass \(1/2\) at each of
\((0,1),(1,0)\). These choices give the two displayed couplings.
In either case each potential outcome is fair, so
\[
 \mathbb P(A=a,Y=b)=\tfrac12\mathbb P(Y(a)=b)=\tfrac14
 \qquad(a,b\in\{0,1\}).
\]
All single-world intervention laws on the two endogenous variables
can be checked exhaustively by their intervention index set:
\begin{enumerate}
\item With no intervention the preceding full table is uniform.
\item With \(A=a\) fixed, \(A\) is constant and \(Y\) is fair,
      since it equals the fair component \(Y(a)\).
\item With \(Y=b\) fixed, \(Y\) is constant and \(A\) remains
      the unchanged fair treatment disturbance.
\item With both fixed, the joint law is the prescribed point mass.
\end{enumerate}
These exhaust all subsets of \(\{A,Y\}\), so the full laws,
and therefore all their marginals, agree in every single world.

Both marginal means are \(1/2\), giving ATE zero. Independence
of the whole potential pair from treatment leaves its distribution
unchanged on \(A=1\), giving ATT zero as well. Benefit is the
cell \((i,j)=(0,1)\), whose masses are zero and \(1/2\).
The factual event \(A=1,Y=1\) has mass \(1/4\) in either
model and forces \(j=1\). In the diagonal coupling this also
forces \(i=1\), whereas in the other coupling it forces \(i=0\).
The stated responder-specific probabilities are therefore zero and
one. These equal experimental laws and unequal targets are explicit
nonidentification witnesses.

For the general bounds, write \(b=\mathbb P(Y(0)=0,Y(1)=1)\).
The two marginal probabilities force all four coupling cells to be
\[
\begin{array}{c|cc}
 &Y(1)=0&Y(1)=1\\\hline
 Y(0)=0&1-p_0-b&b\\
 Y(0)=1&p_0-p_1+b&p_1-b
\end{array}
\]
and these sum to one. Their nonnegativity is precisely
\(b\geq0\), \(b\geq p_1-p_0\), \(b\leq1-p_0\),
and \(b\leq p_1\). This proves the bounds. For
\(p_0,p_1\in[0,1]\), each proposed lower endpoint is no larger
than either proposed upper endpoint, so the interval is nonempty.
Every \(b\) in it gives four nonnegative cells summing to one
with the required margins. Use those cells for \(\kappa\) in
the explicit selector model above. It realizes that benefit
probability with randomized treatment and the specified marginal
potential-outcome laws. Hence the bounds are sharp. When the
interval collapses, this particular target is identified by its
margins; no blanket claim of nonidentification is made for such
degenerate cases.
\end{proof}
\end{annexitem}
